\section{Reference}
\label{sec:reference}

\subsection{Keyboard shortcuts}

\begin{table}[htbp]
  \centering
  \caption{Shortcuts (use \kbd{Ctrl} in place of \kbd{Cmd} on Windows and Linux)}
  \label{tab:shortcuts}
  \begin{tabular}{ll}
    \toprule
    \textbf{Keys} & \textbf{Action} \\
    \midrule
    \kbd{Cmd+Enter}       & Typeset \\
    \kbd{Cmd+S}           & Save now (autosave handles this anyway) \\
    \kbd{Cmd+Shift+P}     & Command palette (lists all shortcuts) \\
    \midrule
    \kbd{Cmd+K}           & Inline AI edit on the selection or paragraph \\
    \kbd{Cmd+Enter}       & \dots{} accept the proposed edit \\
    \kbd{Esc}             & \dots{} discard it, or stop a run \\
    \kbd{Enter}           & \dots{} refine and run again \\
    \midrule
    \kbd{Cmd+F}           & Find \\
    \kbd{Cmd+Alt+F}       & Find and replace \\
    \kbd{Cmd+G}           & Go to line \\
    \kbd{Cmd+/}           & Toggle comment \\
    \kbd{Shift+Alt+Down}  & Duplicate line \\
    \kbd{Alt+Up} / \kbd{Alt+Down} & Move line \\
    \midrule
    \kbd{Cmd+}\texttt{+} / \kbd{Cmd+}\texttt{-} & Zoom the preview in / out \\
    \kbd{Cmd+0}           & Preview back to 100\% \\
    \bottomrule
  \end{tabular}
\end{table}

The preview zoom keys apply while the pointer is over the preview. Both panes
are on screen at once, so the shortcut has to mean whichever one you are
looking at --- with the pointer over the editor, the same keys resize the
editor text instead.

Anything not bound to a key is in the command palette, including toggling the
sidebar, preview and issues panel.

\subsection{Zooming the preview}

Three ways, all doing the same thing:

\begin{itemize}[leftmargin=*]
  \item The \ui{$-$} and \ui{$+$} buttons, which step through a ladder of
        round numbers rather than drifting onto values like 137\%.
  \item Pinch on a trackpad, or \kbd{Ctrl}+scroll, which zooms continuously and
        keeps the point under the pointer still --- so you can zoom into a
        figure by pointing at it.
  \item The percentage itself, which opens presets plus \ui{Fit width} and
        \ui{Fit page}.
\end{itemize}

The two fit options are modes, not one-off calculations: drag the splitter and
the page re-fits to the new pane width. Any manual zoom leaves the mode, since
otherwise the next resize would undo what you just did.

Pages render at your display's full pixel density, so zooming in shows more
detail rather than magnifying a blurry image.

\subsection{Choosing a compiler}

\begin{description}[style=nextline, leftmargin=0pt, labelindent=0pt, labelwidth=0pt, font=\normalfont\ttfamily]
  \item[pdflatex] The default. Fastest, and correct for most documents.
  \item[xelatex] Choose when you need system fonts via \texttt{fontspec}, or
        substantial non-Latin script.
  \item[lualatex] Choose for \texttt{fontspec} plus embedded Lua, or when a
        package requires it.
  \item[latexmk] Not a compiler but a driver: it re-runs the real compiler
        until references settle. The one to pick if the two-press rule of
        Section~\ref{sec:passes} irritates you, or if you need BibTeX run
        automatically.
\end{description}

\subsection{Troubleshooting}

\begin{description}[style=nextline, leftmargin=0pt, labelindent=0pt, labelwidth=0pt, font=\normalfont\itshape]

  \item[The compiler menu is empty.]
  No \LaTeX{} distribution was found on your \texttt{PATH}. Install TeX Live,
  MacTeX or MiKTeX and restart. Verify with \texttt{which pdflatex} in a
  terminal --- if the shell cannot find it either, the problem is the
  installation, not the editor.

  \item[References show as \texttt{??}, the contents page is empty.]
  Press \ui{Typeset} a second time. See Section~\ref{sec:passes}.

  \item[Citations are undefined.]
  The editor does not run BibTeX or Biber. Either use a
  \texttt{thebibliography} environment as this guide does, run
  \texttt{biber}/\texttt{bibtex} once in a terminal, or switch to
  \texttt{latexmk}.

  \item[The wrong file is being typeset.]
  The editor builds the file containing \verb|\documentclass|. If your project
  has several, the shortest path wins. Keep one root document per project
  folder.

  \item[The AI panel says ``whole file'' when you wanted a paragraph.]
  Nothing is selected. Select the text first --- the panel always tells you
  what it is about to send.

  \item[\ui{Apply} is greyed out.]
  Either the model proposed no change, or you switched files after running.
  The panel says which.

  \item[The MCP indicator says ``not running''.]
  Another process holds port 9876. Everything else works.

\end{description}

\subsection{Current limitations}

Stated plainly, so nothing here is a surprise:

\begin{itemize}[leftmargin=*]
  \item One file open at a time --- there are no editor tabs.
  \item One compiler pass per press, with the consequences in
        Section~\ref{sec:passes}.
  \item SyncTeX exists in the backend but is not bound to a click.
  \item Plugins are listed, not executed.
  \item AI context collects \verb|\label|s from the open file only; citation
        keys are gathered across the whole project.
  \item One AI request at a time.
\end{itemize}

\subsection{Further reading}

For \LaTeX{} itself rather than this editor, Lamport~\cite{lamport1994} remains
the clearest introduction, Mittelbach and Goossens~\cite{mittelbach2004} is the
reference to own when a package fights you, and Knuth~\cite{knuth1984} is where
all of it comes from.
